\section{Attention Sink 现象}

\subsection{什么是 Attention Sink？}

\begin{importantbox}{Attention Sink 定义}
在自回归语言模型中，大量 Attention Head 会将异常高的注意力权重分配给序列的\textbf{第一个 token}（或前几个 token），即使这些 token 在语义上并不重要。这种现象称为 \textbf{Attention Sink}。
\end{importantbox}

\subsection{为什么会出现 Attention Sink？}

Softmax 的特性要求注意力权重之和为 1。当模型对某些位置``没什么可注意的''时，它需要把注意力``倾倒''到某个地方。第一个 token 成为了默认的``垃圾桶''：

\begin{warningbox}{Attention Sink 的成因}
\begin{itemize}
  \item Softmax 强制归一化 $\to$ 注意力必须分配到某处
  \item 第一个 token 对所有后续 token 都可见（因果 mask 下）
  \item 模型训练中自然学会将``无意义''的注意力倾倒在此
  \item 并非第一个 token 本身有特殊的语义价值
\end{itemize}
\end{warningbox}

\subsection{Attention Sink 的影响}

\begin{itemize}
  \item \textbf{KV Cache}：在流式推理（sliding window）中，如果丢弃了第一个 token 的 KV，模型输出质量会显著下降
  \item \textbf{模型质量}：注意力资源被浪费在无意义的位置
  \item \textbf{可解释性}：干扰了对模型真实注意力模式的分析
\end{itemize}

\subsection{本章小结}

Attention Sink 是 Softmax 归一化的副产品，第一个 token 成为了注意力的``默认接收器''。这个现象对 KV Cache 管理、模型效率和可解释性都有重要影响。

